\documentclass[10pt,a4paper]{amsart}
\usepackage{amsmath,amssymb,amsthm,graphicx}
\usepackage{verbatim}
\newenvironment{reference}{\comment}{\endcomment}
%\newenvironment{reference}{}{}
\newenvironment{slogan}{\comment}{\endcomment}
\newenvironment{history}{\comment}{\endcomment}
% For commutative diagrams we use Xy-pic
\usepackage[all]{xy}

% We use 2cell for 2-commutative diagrams.
\xyoption{2cell}
% Theorem environments.
%
\theoremstyle{plain}
\newtheorem{theorem}[subsection]{Theorem}
\newtheorem{proposition}[subsection]{Proposition}
\newtheorem{lemma}[subsection]{Lemma}

\theoremstyle{definition}
\newtheorem{definition}[subsection]{Definition}
\newtheorem{example}[subsection]{Example}
\newtheorem{exercise}[subsection]{Exercise}
\newtheorem{situation}[subsection]{Situation}

\theoremstyle{remark}
\newtheorem{remark}[subsection]{Remark}
\newtheorem{remarks}[subsection]{Remarks}

\numberwithin{equation}{subsection}

% Macros
%
\def\lim{\mathop{\mathrm{lim}}\nolimits}
\def\colim{\mathop{\mathrm{colim}}\nolimits}
\def\Spec{\mathop{\mathrm{Spec}}}
\def\Hom{\mathop{\mathrm{Hom}}\nolimits}
\def\Ext{\mathop{\mathrm{Ext}}\nolimits}
\def\SheafHom{\mathop{\mathcal{H}\!\mathit{om}}\nolimits}
\def\SheafExt{\mathop{\mathcal{E}\!\mathit{xt}}\nolimits}
\def\Sch{\mathit{Sch}}
\def\Mor{\mathop{\mathrm{Mor}}\nolimits}
\def\Ob{\mathop{\mathrm{Ob}}\nolimits}
\def\Sh{\mathop{\mathit{Sh}}\nolimits}
\def\NL{\mathop{N\!L}\nolimits}
\def\CH{\mathop{\mathrm{CH}}\nolimits}
\def\proetale{{pro\text{-}\acute{e}tale}}
\def\etale{{\acute{e}tale}}
\def\QCoh{\mathit{QCoh}}
\def\Ker{\mathop{\mathrm{Ker}}}
\def\Im{\mathop{\mathrm{Im}}}
\def\Coker{\mathop{\mathrm{Coker}}}
\def\Coim{\mathop{\mathrm{Coim}}}

% Boxtimes
%
\DeclareMathSymbol{\boxtimes}{\mathbin}{AMSa}{"02}

%
% Macros for moduli stacks/spaces
%
\def\QCohstack{\mathcal{QC}\!\mathit{oh}}
\def\Cohstack{\mathcal{C}\!\mathit{oh}}
\def\Spacesstack{\mathcal{S}\!\mathit{paces}}
\def\Quotfunctor{\mathrm{Quot}}
\def\Hilbfunctor{\mathrm{Hilb}}
\def\Curvesstack{\mathcal{C}\!\mathit{urves}}
\def\Polarizedstack{\mathcal{P}\!\mathit{olarized}}
\def\Complexesstack{\mathcal{C}\!\mathit{omplexes}}
% \Pic is the operator that assigns to X its picard group, usage \Pic(X)
% \Picardstack_{X/B} denotes the Picard stack of X over B
% \Picardfunctor_{X/B} denotes the Picard functor of X over B
\def\Pic{\mathop{\mathrm{Pic}}\nolimits}
\def\Picardstack{\mathcal{P}\!\mathit{ic}}
\def\Picardfunctor{\mathrm{Pic}}
\def\Deformationcategory{\mathcal{D}\!\mathit{ef}}
\usepackage{cleveref}
\crefname{originallemma}{original lemma}{original lemmas}
\Crefname{originallemma}{Original lemma}{Original lemmas}
\makeatletter
\newlabel{decent-spaces-lemma-bounded-fibres}{{03JX}{source}{}{sourcepointer-1}{}}
\newlabel{decent-spaces-lemma-bounded-fibres@cref}{{[originallemma][0][]03JX}{[1][0][]source}{}{}{}}
\makeatother
\crefname{originallemma}{original lemma}{original lemmas}
\Crefname{originallemma}{Original lemma}{Original lemmas}
\makeatletter
\newlabel{lemma-alternating-spectral-sequence}{{0728}{source}{}{sourcepointer-2}{}}
\newlabel{lemma-alternating-spectral-sequence@cref}{{[originallemma][0][]0728}{[1][0][]source}{}{}{}}
\makeatother
\crefname{originallemma}{original lemma}{original lemmas}
\Crefname{originallemma}{Original lemma}{Original lemmas}
\makeatletter
\newlabel{lemma-quasi-coherent-twist}{{072A}{source}{}{sourcepointer-3}{}}
\newlabel{lemma-quasi-coherent-twist@cref}{{[originallemma][0][]072A}{[1][0][]source}{}{}{}}
\makeatother
\crefname{originalproposition}{original proposition}{original propositions}
\Crefname{originalproposition}{Original proposition}{Original propositions}
\makeatletter
\newlabel{proposition-vanishing}{{072B}{source}{}{sourcepointer-4}{}}
\newlabel{proposition-vanishing@cref}{{[originalproposition][0][]072B}{[1][0][]source}{}{}{}}
\makeatother
\crefname{originallemma}{original lemma}{original lemmas}
\Crefname{originallemma}{Original lemma}{Original lemmas}
\makeatletter
\newlabel{spaces-lemma-quotient-finite-separated}{{02Z4}{source}{}{sourcepointer-5}{}}
\newlabel{spaces-lemma-quotient-finite-separated@cref}{{[originallemma][0][]02Z4}{[1][0][]source}{}{}{}}
\makeatother
\crefname{originallemma}{original lemma}{original lemmas}
\Crefname{originallemma}{Original lemma}{Original lemmas}
\makeatletter
\newlabel{spaces-properties-lemma-quasi-compact-affine-cover}{{03H6}{source}{}{sourcepointer-6}{}}
\newlabel{spaces-properties-lemma-quasi-compact-affine-cover@cref}{{[originallemma][0][]03H6}{[1][0][]source}{}{}{}}
\makeatother
\title{Covers and degree bounds for a quasi-compact quasi-separated space, and the resulting vanishing above the maximum of the bounds}
\author{The Stacks Project Authors}
\date{stacks-project:spaces-cohomology.tex:local-snapshot-20261012; GNU FDL 1.2 or later}
\begin{document}
\maketitle
\noindent\textit{Editorial note.} All delivered mathematical text is the author's own, copied from the pinned Stacks Project chapter file spaces-cohomology.tex; no statement or proof step was rewritten, reordered or removed.\par\medskip
\noindent\textit{Editorial note.} Preamble provenance: the pinned chapter file opens with \textbackslash{}input\{preamble\} and carries no preamble text of its own, so the project's single Stacks-wide preamble staged at sources/stacks-spaces-cohomology/preamble.tex is excerpted verbatim as three spans, lines 8-12 (verbatim package and the project's own reference/slogan/history display environments), lines 14-18 (Xy-pic setup) and lines 154-217 (theorem declarations and mathematical macros). The wrapper class and the amsmath/amssymb/amsthm/graphicx base packages are supplied by the assembler. No macro was invented or redefined.\par\medskip
\noindent\textit{Editorial note.} Same-chapter pointers are located at their \textbackslash{}label lines inside the pinned chapter file; cross-chapter pointers are located in the project's own tag list staged in the snapshot (sources/stacks-chow/tags.txt), at the verbatim lines carrying their canonical tags.\par\medskip
\noindent\textit{Allowed standard prerequisites (editorial).}\begin{itemize}
\item An affine etale cover, and the same applied to find affine covers of the configuration quotients. Properties of Spaces, label spaces-properties-lemma-quasi-compact-affine-cover (invoked at proof lines 1406 and 1419). It supplies the affine \$U\$ mapping surjectively and etale onto \$X\$, and later the affine schemes \$V\_p\$ mapping onto the configuration quotients \$U\_p\$.
\item Universal boundedness of the fibres of a morphism from a scheme to a reasonable algebraic space. Decent Spaces, label decent-spaces-lemma-bounded-fibres (invoked at proof lines 1409 and 1422). It produces the integer \$d\$ bounding the degrees of the fibres of \$U \textbackslash{}to X\$ and the integers \$d\_p\$ bounding the degrees of \$V\_p \textbackslash{}to U\_p\$.
\item Quasi-compactness and separatedness of the quotients of the configuration schemes. Spaces, label spaces-lemma-quotient-finite-separated (invoked at proof line 1416). Since \$X\$ is quasi-separated the schemes \$W\_p\$ are quasi-compact and separated, hence the quotients \$U\_p = W\_p/S\_\{p+1\}\$ are too.
\item The alternating Cech spectral sequence, the vanishing on the separated case, and quasi-coherence of the twists. Same chapter spaces-cohomology.tex, labels lemma-alternating-spectral-sequence (invoked at proof line 1434), proposition-vanishing (invoked at proof line 1436) and lemma-quasi-coherent-twist (invoked at proof line 1438). The spectral sequence computes the cohomology of \$X\$ from the twisted cohomology of the \$U\_p\$; the second ingredient makes those groups vanish above \$d\_p\$; the third keeps the twisted sheaves quasi-coherent so that the vanishing applies.
\end{itemize}\medskip
\ifdefined\hypertarget\hypertarget{sourcepointer-1}{}\fi
\noindent\textit{Source reference.} Original lemma 03JX: sources/stacks-chow/tags.txt, source lines 4655--4655. Navigation pointer for \textbackslash{}ref\{decent-spaces-lemma-bounded-fibres\} used inside the counted span (line 1409); the argument that consumes it is retained in the span.\par\medskip
\ifdefined\hypertarget\hypertarget{sourcepointer-2}{}\fi
\noindent\textit{Source reference.} Original lemma 0728: sources/stacks-spaces-cohomology/pinned-original-source.tex, source lines 1217--1217. Navigation pointer for \textbackslash{}ref\{lemma-alternating-spectral-sequence\} used inside the counted span (line 1395); the argument that consumes it is retained in the span.\par\medskip
\ifdefined\hypertarget\hypertarget{sourcepointer-3}{}\fi
\noindent\textit{Source reference.} Original lemma 072A: sources/stacks-spaces-cohomology/pinned-original-source.tex, source lines 1305--1305. Navigation pointer for \textbackslash{}ref\{lemma-quasi-coherent-twist\} used inside the counted span (line 1438); the argument that consumes it is retained in the span.\par\medskip
\ifdefined\hypertarget\hypertarget{sourcepointer-4}{}\fi
\noindent\textit{Source reference.} Original proposition 072B: sources/stacks-spaces-cohomology/pinned-original-source.tex, source lines 1333--1333. Navigation pointer for \textbackslash{}ref\{proposition-vanishing\} used inside the counted span (line 1436); the argument that consumes it is retained in the span.\par\medskip
\ifdefined\hypertarget\hypertarget{sourcepointer-5}{}\fi
\noindent\textit{Source reference.} Original lemma 02Z4: sources/stacks-chow/tags.txt, source lines 3899--3899. Navigation pointer for \textbackslash{}ref\{spaces-lemma-quotient-finite-separated\} used inside the counted span (line 1416); the argument that consumes it is retained in the span.\par\medskip
\ifdefined\hypertarget\hypertarget{sourcepointer-6}{}\fi
\noindent\textit{Source reference.} Original lemma 03H6: sources/stacks-chow/tags.txt, source lines 4555--4555. Navigation pointer for \textbackslash{}ref\{spaces-properties-lemma-quasi-compact-affine-cover\} used inside the counted span (line 1407); the argument that consumes it is retained in the span.\par\medskip
\begin{lemma}
\label{lemma-vanishing-quasi-separated}
Let $S$ be a scheme. Let $X$ be an algebraic space over $S$.
Assume $X$ is quasi-compact and quasi-separated.
Then we can choose
\begin{enumerate}
\item an affine scheme $U$,
\item a surjective \'etale morphism $f : U \to X$,
\item an integer $d$ bounding the degrees of the fibres of $U \to X$,
\item for every $p = 0, 1, \ldots, d$ a surjective \'etale morphism
$V_p \to U_p$ from an affine scheme $V_p$ where $U_p$ is as in
Lemma \ref{lemma-alternating-spectral-sequence}, and
\item an integer $d_p$ bounding the degree of the fibres of $V_p \to U_p$.
\end{enumerate}
Moreover, whenever we have (1) -- (5), then for any quasi-coherent
$\mathcal{O}_X$-module $\mathcal{F}$ we have $H^q(X, \mathcal{F}) = 0$ for
$q \geq \max(d_p + p)$.
\end{lemma}

\begin{proof}
Since $X$ is quasi-compact we can find a surjective \'etale morphism
$U \to X$ with $U$ affine, see
Properties of Spaces, Lemma
\ref{spaces-properties-lemma-quasi-compact-affine-cover}.
By
Decent Spaces, Lemma \ref{decent-spaces-lemma-bounded-fibres}
the fibres of $f$ are universally bounded, hence we can find $d$.
We have $U_p = W_p/S_{p + 1}$ and $W_p \subset U \times_X \ldots \times_X U$
is open and closed. Since $X$ is quasi-separated the schemes
$W_p$ are quasi-compact, hence $U_p$ is quasi-compact.
Since $U$ is separated, the schemes $W_p$ are separated, hence
$U_p$ is separated by (the absolute version of)
Spaces, Lemma \ref{spaces-lemma-quotient-finite-separated}.
By
Properties of Spaces, Lemma
\ref{spaces-properties-lemma-quasi-compact-affine-cover}
we can find the morphisms $V_p \to W_p$.
By
Decent Spaces, Lemma \ref{decent-spaces-lemma-bounded-fibres}
we can find the integers $d_p$.

\medskip\noindent
At this point the proof uses the spectral sequence
$$
E_1^{p, q} =
H^q(U_p, \mathcal{F}|_{U_p} \otimes_\mathbf{Z} \underline{\mathbf{Z}}(\chi_p))
\Rightarrow
H^{p + q}(X, \mathcal{F})
$$
see
Lemma \ref{lemma-alternating-spectral-sequence}.
By definition of the integer $d$ we see that $U_p = 0$ for $p \geq d$.
By Proposition \ref{proposition-vanishing}
and
Lemma \ref{lemma-quasi-coherent-twist}
we see that
$H^q(U_p,
\mathcal{F}|_{U_p} \otimes_\mathbf{Z} \underline{\mathbf{Z}}(\chi_p))$
is zero for $q \geq d_p$ for $p = 0, \ldots, d$.
Whence the lemma.
\end{proof}

\end{document}
